\begin{helper}
Work in the same single exposure-history cell setting as
\(\noderef{FixedSetHistoryCellAdaptiveProductBound}\), with the same
cylinder description, exact terminal product law, prefix-safety hypothesis,
regular open-pair lower bound, and branch-indexed adaptive exceptional
candidate package.  Fix natural numbers \(u_R,u_B\).  Then the conditional
probability, inside the cell, of final independence and regularity together
with exactly \(u_R\) present red staged-open terminal coordinates and exactly
\(u_B\) present blue staged-open terminal coordinates is at most
\[
\exp\!\left(
8\sqrt{\varepsilon_1}\,pk^2
-p\bigl(f(\ell_R,\ell_B)-10\varepsilon_1k^2\bigr)\right).
\]
\end{helper}
\begin{proof}
Put \(a=f(\ell_R,\ell_B)-10\varepsilon_1k^2\).  The regular open-pair lower
bound gives \(a\le |S(\omega)|\) for every successful regular sample, where
\(S(\omega)\) is its staged-open terminal set.  The finite cell partition,
per-cell fixed-cylinder charge, and total branch/transcript weight sum used
below are packaged in
\(\noderef{FixedSetHistoryCellFixedCountSummationBridge}\).  Let
\[
P_R(\omega)=\{e\in S(\omega):e\hbox{ is red and present}\},
\qquad
P_B(\omega)=\{e\in S(\omega):e\hbox{ is blue and present}\}.
\]
On the fixed-count event, \(|P_R(\omega)|=u_R\) and
\(|P_B(\omega)|=u_B\).  Split into the two cases \(u_R\ge u_B\) and
\(u_B\ge u_R\).

Assume first that \(u_R\ge u_B\).  We expose the blue side first, but every
probability factor is charged on a fixed finite cylinder inside the original
terminal product cell.  Fix a possible blue present support
\[
B\subseteq U_{\mathrm{blue}},\qquad |B|=u_B .
\]
The blue staged-open coordinates are projected pairs from the \(k\)-set, so
the ambient blue pool has size \(N_B\le k^2\), and there are at most
\(\binom{N_B}{u_B}\) possible choices for \(B\).
The fixed-cylinder charge
\(\noderef{FixedSetHistoryCellFixedCylinderCharge}\), applied with present
set \(B\) and empty absent set, gives a factor \(p^{u_B}\).  All terminal
coordinates not mentioned by this fixed requirement are summed out in the
same terminal product cell; no complete blue branch has been conditioned on
yet.

Now refine the fixed blue-support event into complete blue terminal traces,
but do not normalize separately inside each trace.  We keep the calculation
as one finite weighted sum over the disjoint complete traces and the later
prefix transcripts; the trace cover and disjointness are supplied by
\(\noderef{FixedSetHistoryCellAdaptiveBranchPartition}\).  Choose a
representative \(\sigma\) for a nonempty complete blue trace.  On this fixed
trace, the branch-indexed candidate hypothesis supplies a red exceptional set
\(E_R(\sigma)\subseteq U\) with
\[
|E_R(\sigma)|\le 3\varepsilon_1 k^2
\]
because \(\noderef{ClosedPairsBound}( |E_R(\sigma)|,3\varepsilon_1,k)\)
means exactly this inequality.  The set \(E_R(\sigma)\) is fixed before any
red terminal value is charged.  Every present red staged-open coordinate in a
regular independent completion extending the trace lies in \(E_R(\sigma)\),
so for this trace there are at most \(\binom{N_R}{u_R}\) red supports
\[
R\subseteq E_R(\sigma),\qquad |R|=u_R,
\]
with \(N_R\le 3\varepsilon_1k^2\).

Fix such a complete blue trace and red support \(R\).  For a successful
completion with these fixed data, the lower bound \(a\le |S(\omega)|\) and
the fixed-count identities imply that at least
\[
\max(0,a-u_R-u_B)
\]
coordinates of \(S(\omega)\setminus(R\cup B)\) are absent.  These absent
coordinates are not charged as a branch-dependent set after the fact.  Scan
the terminal coordinates in the fixed order.  When the scan reaches a
coordinate \(e\), the prefix-safety hypothesis says that staged-openness,
touchedness, same-color closedness, and recordedness of \(e\) are determined
from the history and earlier terminal values.  Therefore the decision that
\(e\) is a staged-open coordinate outside the already fixed supports
\(R\cup B\) is made before the Bernoulli value of \(e\) is read.

The canonical selection
\(\noderef{FixedSetHistoryCellCanonicalAbsenceSelection}\) is implemented as
the stopping rule that selects the first required coordinates declared by
this prefix test.  It does not search for coordinates after observing that
they are absent.  Partition the successful completions by triples
\[
(\sigma,R,\tau),
\]
where \(\sigma\) is the complete blue trace, \(R\) is the red support, and
\(\tau\) is the prefix transcript up to the selected coordinates.  In a fixed
cell \(G_{\sigma,R,\tau}\), all terminal values mentioned by the blue trace,
all red coordinates in \(R\), all earlier transcript values, and the selected
absent coordinates \(Z(\tau)\) are fixed.  If a selected absent coordinate is
blue, it is the same fixed absence already appearing in the complete blue
trace; it is not a second requirement.  Thus every coordinate contributes at
most one Bernoulli factor in the combined fixed cylinder.

Apply \(\noderef{FixedSetHistoryCellFixedCylinderCharge}\) to each fixed
cylinder \(G_{\sigma,R,\tau}\).  The cells cover the successful event for
the fixed blue support \(B\), and their fixed-cylinder weights are summed by
\(\noderef{FixedSetHistoryCellBranchTranscriptSummation}\).  The weight sum
is computed in the scan order.  For coordinates not forced by the support or
by the selected-absence rule, the two possible terminal values are both
included and contribute \(p+(1-p)=1\).  The support coordinates in \(B\)
contribute \(p^{u_B}\), the support coordinates in \(R\) contribute
\(p^{u_R}\), and every selected absent coordinate contributes one factor
\((1-p)\), whether that absence is recorded in the blue trace or in the later
red transcript.  Since
\[
|Z|\ge\max(0,a-u_R-u_B)
\]
on every successful transcript, and \(0\le1-p\le1\), each transcript weight
is bounded by replacing \(|Z|\) with \(\max(0,a-u_R-u_B)\).  Therefore the
total contribution for this fixed blue support \(B\), after summing all
complete blue traces, all red supports, and all prefix transcripts, is at
most
\[
\binom{N_R}{u_R}p^{u_R}p^{u_B}
(1-p)^{\max(0,a-u_R-u_B)}.
\]
There is no multiplier for the number of complete blue traces: their
probabilities are part of the fixed-cylinder weight sum, and all terminal
coordinates of the traces not fixed by \(B\) or by the selected absences sum
out to \(1\).

Thus the fixed-count probability in this case is bounded by
\[
\binom{N_R}{u_R}p^{u_R}
\binom{N_B}{u_B}p^{u_B}
(1-p)^{\max(0,a-u_R-u_B)}
\]
with \(N_R\le3\varepsilon_1k^2\) and \(N_B\le k^2\).  Applying
\(\noderef{FixedSetHistoryCellBinomialExponentialEstimate}\) gives the
claimed exponential bound.  In the case \(u_B\ge u_R\), interchange the
colors: first pay for the red present support in the fixed terminal product
cell, branch on complete red terminal traces, use the fixed set
\(E_B(\sigma)\) to control the later blue support, and run the same
prefix-safe transcript partition before charging any absence factor.  The
displayed bound is identical, so the fixed-count estimate holds for every
ordered pair \((u_R,u_B)\).
\end{proof}
